\documentclass[12pt]{article}

%% Basic document formatting
\usepackage{amsmath, amsthm, amssymb, amsfonts}
\usepackage{mathtools}
\usepackage{xspace}
\usepackage{thmtools}
\usepackage{graphicx}
\usepackage{tikz}
\usepackage{ytableau}
\usepackage{blkarray}
\usepackage{hyperref}
\usepackage{setspace}
\usepackage{array}
\usepackage{fancyhdr}
\usepackage{titling}
\usepackage[left=0.4in,right=0.4in,top=1in,bottom=1in]{geometry}
\usepackage{float}
\usepackage{cancel}
\usepackage{tabularx}
\usepackage[utf8]{inputenc}
\usepackage[english]{babel}
\usepackage{framed}
\usepackage[dvipsnames]{xcolor}
\usepackage{environ}
\usepackage{tcolorbox}
\tcbuselibrary{theorems,skins,breakable}

\usepackage{enumitem}
\setlist[enumerate]{leftmargin=*}
\setlist[enumerate,1]{labelindent=\parindent}
\setlist[enumerate,2]{labelindent=0pt}

% Blackboard Bold
\newcommand{\Z}{\mathbb{Z}}                 % integers
\newcommand{\Zpl}{\mathbb{Z}_{+}}           % positive integers
\newcommand{\Znn}{\mathbb{Z}_{\geq 0}}      % nonnegative integers. 
\newcommand{\Q}{\mathbb{Q}}                 % rationals
\newcommand{\Qpl}{\mathbb{Q}_{+}}           % positive Rationals
\newcommand{\R}{\mathbb{R}}                 % reals
\newcommand{\Rpl}{\mathbb{R}_{+}}           % positive Reals
\newcommand{\C}{\mathbb{C}}                 % complex numbers
\newcommand{\F}{\mathbb{F}}                 % field

% Words
\newcommand{\st}{\text{ such that }}
\newcommand{\wrt}{\text{ with respect to }}
\newcommand{\with}{\text{ with }}
\newcommand{\ie}{\text{, i.e. }}
\newcommand*{\ditto}{\text{---}\texttt{"}\text{---}}

% Operators
\newcommand{\abs}[1]{\left|#1\right|}                   % absolute value
\newcommand{\norm}[1]{\abs{\abs{#1}}}
\newcommand{\floor}[1]{\left\lfloor #1 \right\rfloor}   % floor
\newcommand{\ceil}[1]{\left\lceil #1 \right\rceil}      % ceiling
\newcommand{\powset}[1]{\mathscr{P}(#1)}

% Category Theory 
\renewcommand{\hom}{\text{Hom}}
\newcommand{\homspace}[3]{\hom_{#1}(#2, #3)}
\newcommand{\coproduct}[9]{
  \begin{tikzcd}[ampersand replacement=\&]
      \& \& #1 \arrow[dll, bend right, "#6"] \arrow[dl, "#5"] \\
   #4 \& #3 \arrow[l, "#9"] \\
      \& \& #2 \arrow[ull, bend left, "#8"] \arrow[ul, "#7"]
  \end{tikzcd}
}

\newcommand{\product}[9]{ 
  \begin{tikzcd}[ampersand replacement=\&]
      \& \& #3 \\
      #1 \arrow[r, "#7"] \arrow[urr, bend left, "#5"] \arrow[drr, bend right, "#6"] \& #2 \arrow[ur, "#8"] \arrow[dr, "#9"] \\ 
      \& \& #4
  \end{tikzcd}
}


% Algebra 
\newcommand{\Syl}[2]{\text{Syl}_{#1}{(#2)}}           % set of Sylow-p subgroups 
\newcommand{\<}{\langle}                % \<x\>, subgroup generated by x 
\renewcommand{\>}{\rangle}
\renewcommand{\index}[2]{[#1 : #2]}
\newcommand{\id}{\text{id}}             % identity element
\newcommand{\lhdneq}{%
  \mathrel{\ooalign{$\lneq$\cr\raise.22ex\hbox{$\lhd$}\cr}}}
\newcommand{\order}[1]{\text{o}(#1)}    % order of an element
\newcommand{\spec}{\operatorname{Spec}}
\newcommand{\rad}{\operatorname{rad}}   % radical of an ideal 
\newcommand{\sym}[1]{\mathfrak{S}_{#1}}
\newcommand{\aut}[1]{\text{Aut}(#1)} 
\newcommand{\gal}[2]{\text{Gal}(#1 / #2)} 
\newcommand{\inn}[1]{\text{Inn}(#1)} 
\let\oldcong\cong
\let\oldequiv\equiv
\renewcommand{\cong}{\oldequiv}
\renewcommand{\equiv}{\oldcong}

\newcommand{\triangleleftneq}{\mathrel{\ooalign{%
  $\trianglelefteq$\cr
  \hidewidth\raise0.06ex\hbox{$\not\mkern4mu$}\hidewidth\cr
}}}

\makeatletter % cycle
\newcommand{\cyc}[1]{(\mathbf{\cyc@process#1\relax})}
\def\cyc@process#1#2\relax{%
	#1%
	\ifx\relax#2\relax
	\else
		\,\cyc@process#2\relax
	\fi
}
\makeatother

\newcommand{\GL}[2]{\text{GL}_{#1}(#2)}                             % general linear group
\newcommand{\SL}[2]{\text{SL}_{#1}(#2)}                             % special linear group
\newcommand{\gl}[2]{\mathfrak{gl}_{#1}(#2)}
\renewcommand{\sl}[2]{\mathfrak{sl}_{#1}(#2)}
\newcommand{\so}[2]{\mathfrak{so}_{#1}(#2)}
\newcommand{\su}[1]{\mathfrak{su}(#1)}

% Linear Algebra
\newcommand{\bmat}[1]{\begin{bmatrix}#1\end{bmatrix}}   % bracketed matrix
\newcommand{\rank}{\operatorname{rank}}                 % rank
\newcommand{\nullity}{\operatorname{nullity}}           % nullity
\newcommand{\tr}{\operatorname{tr}}

% Combinatorics
\newcommand{\qnum}[1]{\left[#1\right]_q}                % q-analog
\newcommand{\qfact}[1]{\left[#1\right]!_q}              % q-analog factorial 
\newcommand{\qbinom}[2]{\binom{#1}{#2}_q}               % Gaussian binomial coefficient

% Topology/Analysis
\newcommand{\ball}[2]{\text{B}_{#1}(#2)}  % B_r(x): open r-balls around x
\newcommand{\diam}{\text{diam}}           % diamter of a set in metric space 
\newcommand{\lowint}[2]{\underline{\int_{#1}^{#2}}}
\newcommand{\upint}[2]{\overline{\int}_{#1}^{#2}}

% Misc Notation
\newcommand{\oneton}[1]{\{1, \ldots, #1\}}
\newcommand{\defeq}{\vcentcolon=}   % :=
\newcommand{\eqdef}{=\vcentcolon}   % =: 
\renewcommand{\bf}[1]{\textbf{#1}}
\newcommand{\im}{\operatorname{im}}
\newcommand{\fnrest}[2]{\left.#1\right|_{#2}}
\newcommand{\isom}{\overset{\sim}{\rightarrow}} 


\renewcommand{\thesection}{\S \arabic{section}}
\renewcommand{\thesubsection}{\arabic{subsection}.}
\renewcommand{\thesubsubsection}{(\alph{subsubsection})} 

\newtheorem{claim}{Claim}
\newtheorem*{lemma}{Lemma}

%% Headers & title setup
\newcommand{\course}{\textit{Topology}, Munkres} 
\newcommand{\myname}{Jay Ser}
\setlength{\headheight}{14.5pt}
\pagestyle{fancy}
\fancyhf{}
\renewcommand{\headrulewidth}{0.4pt}
\lhead{\course}
\rhead{\myname}
\cfoot{\thepage}
\setlength{\droptitle}{-4em} 
\title{\course\ - Chapter 3}
\author{\myname}
\date{March 24, 2026 $\sim$ \today}

\begin{document}
\maketitle
\thispagestyle{fancy}

%------------------------------------------------------------------------------%
\setcounter{section}{22}
\section{Connected Spaces}
\setcounter{subsection}{5}
\subsection{}
Suppose $C \cap \partial A = \emptyset$. 
Then I claim 
\begin{equation} \label{eq:23_6_decomp}
  C = (C \cap (X \setminus \bar{A})) \cup (C \cap A^{\circ})
\end{equation}
gives a separation of $C$. 
It is clear that $X \setminus \bar{A}$ and $A^{\circ}$ are disjoint open sets in $X$, so the only thing to check is that $C \cap (X \setminus \bar{A})$ and $C \cap A^{\circ}$ are nonempty and that their union is indeed equal to $C$. 
Since $C \cap A \neq \emptyset$, take $x \in C \cap A$. 
$C \cap \partial{A} = \emptyset \implies$ there is a neighborhood $U_x$ of $x$ such that $U_x \cap A = \emptyset$ or $U_x \cap (X \setminus A) = \emptyset$. 
But $x \in A$ means $U_x \cap U_x \cap A \neq \emptyset$, so $U_x \cap (X \setminus A) = \emptyset \implies U_x \subset A$. 
This shows $x \in C \cap A^{\circ}$, hence $C \cap A^{\circ} \neq \emptyset$. 
Symmetric argument shows that if $y \in C \cap (X \setminus A)$, there is a neighborhood $U_y$ of $y$ such that $U_y \cap A = \emptyset \implies y \notin \bar{A} \implies C \cap (X \setminus \bar{A}) \neq \emptyset$. 

Finally, to show that \eqref{eq:23_6_decomp} is indeed an equality between sets, take $x \in C \setminus (C \cap A^{\circ})$. 
Then for every neighborhood $U_x$ of $x$ in $C$, $U_x \cap (X \setminus A) \neq \emptyset$. 
Since $C \cap \partial{A} = \emptyset$, $\exists U_x \st U_x \cap A = \emptyset \implies x \notin \bar{A} \implies x \in C \cap (X \setminus \bar{A})$, as was to be shown. 
Hence \eqref{eq:23_6_decomp} gives a separation of $C$, $C$ is not connected, which contradicts the assumption that $C$ is connected. 

\setcounter{subsection}{11}
\subsection{} 
Since $A, B \subset X \setminus Y$, $A \cap Y = \emptyset$ and $B \cap Y = \emptyset$. 
Suppose $C$ and $D$ separate $Y \cup B$. 
Since $Y \subset Y \cup B$ is connected, $Y \subset C$ or $Y \subset D$. 
Without loss of generality, assume $Y \subset D$.
Then $A \cup Y = C \cup D$ and $C \cap D = \emptyset$ says $C \subset A$. 
This is enough to show that 
$$(B \cup D) \cup C = X$$
gives a separation of $X$. 
That $B \cup D$ and $C$ are disjoint and that neither $B \cup D$ nor $C$ are empty are obvious. 
It must be shown that 
$$\overline{B \cup D} \cap C = \emptyset \text{ and } (B \cup D) \cap \overline{C} = \emptyset.$$
Recall $\overline{B \cup D} = \overline{B} \cup \overline{D}$. 
Since $C \subset A$, $\overline{B} \cap C \subset \overline{B} \cap A = \emptyset$, the latter equality due to the fact that $A$ and $B$ separate $X \setminus Y$. 
$\overline{C} \cap D = \emptyset$ since $C$ and $D$ separate $A \cup Y$. 
This shows $\overline{B \cup D} \cap C = \emptyset$. 
Next, $B \cap \overline{C} = \emptyset = D \cap \overline{C}$ for similar reasons, hence 
$(B \cup D) \cap \overline{C} = \emptyset$. 
This shows $B \cup D$ and $C$ separate $X$, contradicting the assumption that $X$ is connected. 

\end{document}
